\section{METODOLOGIA}

\subsection{Classificação da Pesquisa}

Esta pesquisa classifica-se como \textbf{aplicada} quanto à sua natureza, pois visa resolver um problema prático específico --- a análise automatizada de falhas em IaC. Quanto aos objetivos, é \textbf{exploratória e descritiva}, uma vez que investiga a aplicabilidade de LLMs em um domínio pouco explorado e descreve quantitativamente o desempenho de diferentes modelos. Quanto à abordagem, adota uma perspectiva \textbf{quantitativa}, utilizando métricas numéricas para comparação. Os procedimentos técnicos incluem \textbf{experimentação controlada} em ambiente laboratorial.

\subsection{Ambiente Experimental}

O ambiente experimental consiste em:
\begin{itemize}
    \item \textbf{Hardware:} Computador pessoal com processador multi-core, 16\,GB de RAM e GPU integrada, executando sistema operacional Linux (Ubuntu).
    \item \textbf{Motor de IA:} Ollama instalado localmente, servindo modelos via API REST na porta 11434.
    \item \textbf{Modelos avaliados:} CodeLlama 7B (Meta), DeepSeek-Coder 6.7B (DeepSeek) e Qwen2.5-Coder 7B (Alibaba).
    \item \textbf{Ferramenta de IaC:} Terraform com provedor Docker para simulação de cenários locais.
    \item \textbf{Linguagem de desenvolvimento:} Python 3.x com bibliotecas Streamlit, Plotly, Pandas e FPDF.
\end{itemize}

\subsection{Cenários de Falhas}

Foram catalogados cenários controlados de falhas em Terraform, organizados em categorias:
\begin{itemize}
    \item \textbf{Erros de sintaxe:} Ausência de chaves, aspas não fechadas, blocos incompletos.
    \item \textbf{Erros de dependência:} Dependências cíclicas entre recursos, referências a recursos inexistentes.
    \item \textbf{Conflitos de recursos:} Portas duplicadas, nomes de containers em conflito.
    \item \textbf{Erros de provedor:} Imagens Docker inexistentes, configurações inválidas de rede.
    \item \textbf{Erros de variáveis:} Variáveis não declaradas, tipos incompatíveis, valores padrão ausentes.
\end{itemize}

Cada cenário consiste em um arquivo \texttt{main.tf} com um erro proposital e um arquivo de metadados descrevendo o tipo de falha esperada.

\subsection{Pipeline de Análise}

O pipeline automatizado opera da seguinte forma:
\begin{enumerate}
    \item O script de automação (\texttt{loop\_runner.py}) seleciona um cenário de falha e um modelo de IA em esquema \textit{round-robin}.
    \item O Terraform é executado no cenário selecionado (\texttt{terraform init} seguido de \texttt{terraform plan/apply}).
    \item O log de erro resultante é capturado automaticamente via \texttt{subprocess}.
    \item O script lê o conteúdo do arquivo \texttt{main.tf} correspondente.
    \item Um prompt estruturado é montado combinando: (a) o papel do modelo (engenheiro DevOps sênior), (b) o código-fonte completo, (c) o log de erro integral, e (d) as instruções de formato de resposta.
    \item O prompt é enviado à API local do Ollama (\texttt{localhost:11434/api/generate}).
    \item A resposta do modelo, o tempo de processamento e a estimativa de \textit{tokens} são registrados em arquivo CSV e, opcionalmente, em banco PostgreSQL.
    \item O ciclo repete-se para cada combinação modelo/cenário pelo número de iterações configurado.
\end{enumerate}

\subsection{Métricas de Avaliação}

Os modelos são avaliados segundo as seguintes métricas:
\begin{itemize}
    \item \textbf{Assertividade:} Capacidade de identificar corretamente a causa raiz da falha (avaliação qualitativa por inspeção do pesquisador).
    \item \textbf{Tempo de resposta:} Tempo total de processamento da IA, em segundos, desde o envio do prompt até a resposta completa.
    \item \textbf{Volume de \textit{tokens}:} Quantidade estimada de \textit{tokens} na resposta gerada.
    \item \textbf{Throughput (\textit{tokens}/segundo):} Razão entre o volume de \textit{tokens} e o tempo de resposta, indicando a eficiência computacional.
    \item \textbf{Aderência à documentação:} Grau em que as correções sugeridas seguem as melhores práticas da documentação oficial do Terraform.
\end{itemize}

\subsection{Análise dos Resultados}

Os dados coletados são analisados por meio de:
\begin{itemize}
    \item Estatísticas descritivas (média, mediana, moda, desvio-padrão, mínimo e máximo) globais e por modelo.
    \item Filtragem automática de \textit{outliers} via método IQR (Intervalo Interquartil) e limite no percentil 97.
    \item Visualizações gráficas: \textit{box plots}, gráficos \textit{violin}, histogramas de distribuição, \textit{heatmaps}, gráficos de dispersão e gráficos radar comparativos.
    \item Exportação dos resultados em formato CSV e PDF para inclusão em relatórios acadêmicos.
\end{itemize}
